\lecture{6}{Thu 08 Sep 2022 12:00}{}

General Principle:
\begin{theorem}
	If $u$ is harmonic, $f$ is analytic, then $u(f)$ is harmonic.
\end{theorem}
\begin{proof}
	\begin{enumerate}
		\item Show directly that $\Delta u =0$ and CR condition imply \[
			\Delta \left( u(f) \right)  = 0
		.\] 
	\item Locally, $u = \operatorname{Re}  F$ for some analytic $F$, then $u(f) = \operatorname{Re} F(f)$. $F(f)$ is analytic, since $F$, $f$ are.
	\end{enumerate}
\end{proof}

\begin{theorem}
	If $\Delta u =0$ and $\Delta v =0$, then $\Delta\left( c_1 u + c_2 v \right) =0$, where $c_1, c_2$ are constants.
\end{theorem}

\begin{example}
	$\mathcal{D} = \{x+iy:0<y<1\} $. We have $u: \overline{\mathcal{D}} \to \mathbb{R}$. Find $u(x,y)$ such that $u(x,0) = 3$ and $u(x,1) = -1$.

	Look for $u(x,y) = a y + b$. This satisfies the boundary conditions: For $y = 0$, $u(x,y) = b = 3$. For  $y = 1$, $u(x,y) = a+b = -1$, so $a = -4$.
\end{example}

\begin{example}
	See problem in the figure.
\begin{figure}[H]
    \centering
    \incfig{dirich3}
    \caption{dirich3}
    \label{fig:dirich3}
\end{figure}
	We construct a function \[
		\varphi(z) = (1+i)z
	\] 
	that is a linear transformation equating this problem and the last problem. We simply apply transformation to the previous solution:
	 \[
		u(x,y) = 2 \operatorname{Im} \varphi(z)
	.\] 
\end{example}

\begin{definition}
	A \textit{curve} is a continuous map $\gamma : [a,b] \to  \mathbb{C}$.
\end{definition}
\begin{figure}[ht]
    \centering
    \incfig{curve}
    \caption{curve}
    \label{fig:curve}
\end{figure}


\begin{definition}[piecewise smooth]
	Exists $a\le t_1<t_2<\ldots<t_n = b$ such that for each $i$, exists $\gamma': [t_i, t_{i+1}] \to \mathbb{C}$ such that $\gamma'(t) \neq 0$.
\end{definition}



Let $\gamma$ be a curve, $f(z)$ be a analytic function. Now \[
	\frac{\partial }{\partial t} f(\gamma(t)) = f'(\gamma(t))\gamma(t)
.\] 
Geometric meaning: tangent vector $\gamma'(t)$ of $\gamma(t)$ is multiplied by complex number $f'(\gamma(t))$.

\begin{definition}
We define the \textit{angle} between two intersecting curves at a point to be the angle between their tangent vectors at this point.
\end{definition}

\begin{theorem}
	Analytic functions preserve angles between curves where $f'(z) \neq  0$. Hence analytic functions are \textbf{conformal}.
\end{theorem}

\begin{nonexample}
	\begin{itemize}
		\item 
	$f(z) = z^2$ increases angle by factor of $2$ at $0$.
\item $\overline{z}$ is not conformal.
	\end{itemize}
\end{nonexample}

\begin{remark}
	In this course we only consider \textit{real harmonic functions}.

\end{remark}
\section{Elementary Functions}

\subsection{Polynomials}

\begin{definition}
	A polynomial is of the form \[
		f(z) = a_d z^d + a_{d-1} z^{d-1} + \ldots+a_0
	.\] where $a_j \in \mathbb{C}$ and $a_d \neq  0$.
	Define the \textit{degree} of polynomial as \[
		\operatorname{deg} f = d
	.\] 
	Sometimes it's convenient to let $\operatorname{deg} (0) = -\infty $.
\end{definition}

Note that $\operatorname{deg} (fg) = \operatorname{deg} f \operatorname{deg} g$.

\begin{theorem}[Division algorithm for polynomials]
	If $f,g$ are polymials, then $f = q g + r$ where $q$ and $r$ are polynomials and $\operatorname{deg}  r < \operatorname{deg} g$. If $r = 0$, we say $f$ is \textit{divisible} by $g$.
\end{theorem}

\begin{definition}[root]
	If $P(z_0) = 0$ and $P\neq 0$, then we say $z_0$ is a \textit{root} of $P$.
\end{definition}

\begin{proposition}
	If $P$ has a root $z_0$, then $(z-z_0) \mid P$.
\end{proposition}
\begin{proof}
Apply the division algorithm, $P = Q(z-z_0) + r(z)$, $\operatorname{deg} r < 1 \implies r $ is constant. Plug $z = z_0$, now $r = 0$. 
\end{proof}

$P = (z - z_0)^m Q^*, Q^*(z_0) \neq  0$ is called the \textit{canonical representation} of $P$ at $z_0$. $m$ is the \textit{multiplicity} of the root $z_0$ of $P$.

If $m = 1$ then it's called a \textit{simple root}.

\begin{definition}
	The \textit{Maclaurin Formula} is \[
		f(z) = \sum_{n=0}^{\infty} \frac{1}{n!}f^{(n)}(0)
	.\] 
\end{definition}

\begin{proposition}
	If $a \in \mathbb{C}$ and $P - a$ is a polynomial, then $P(z-a)$ is a polynomial.
\end{proposition}
Now let $b_n = P^{(n)}(a) / n!$ and we have the \textit{Taylor formula}: \[
	f(z) = \sum_{n=0}^{\infty} b_n (z-a)^n
.\] 

\begin{theorem}[Fundamental Theorem of Algebra]
	Every (complex) polynomial of degree $\ge 1$ has a (complex) root.
\end{theorem}
\begin{corollary}
	Every polynomial completely factors: \[
		P(z) = a(z-z_1)(z-z_2)\ldots(z-z_d)
	\] where $a, z_1, \ldots, z_d$ are complex. 
	We may write it in distinct roots to power of multiplicity:
	\[
		P(z) = a(z-z_1)^{m_1}\ldots(z-z_k)^{m_k}
	\] where $z_1, \ldots, z_k$ are distinct, and $m_1 + \ldots + m_k = d$.
\end{corollary}

\subsection{Rational Functions}

\begin{definition}[Rational Function]
	A \textit{rational function} is a ratio of polynomials: \[
		R = \frac{P}{Q}
	.\] 
	If $P$ and $Q$ have no common roots, then $\frac{P}{Q}$ is a \textit{reduced representation} of $R$.
\end{definition}

\begin{remark}
	to find out if $P$ and $Q$ have common roots is simple (GCD algorithm). On the contrary, to factorize a polynomial is difficult. The same is for integers.
\end{remark}

If $Q(b) = 0$, then $\lim_{z \to b} \frac{P}{Q} = \infty $. We can view $R$ as a map $\mathbb{C} \to \overline{\mathbb{C}} = \mathbb{C}\cup \{\infty \}  $.

$b$ is called a  \textit{pole} of multiplicity $m$ of $R$ if $b$ is a root of multiplicity $m$ of $Q$.

We further extend rational functions to $R:\overline{\mathbb{C}} \to \overline{\mathbb{C}} $.

Define \[
	R(\infty ) = \lim_{z  \to 0} R(\frac{1}{z})
.\] This can be finite or infinite.

\begin{example}
	$f(z) = \frac{1}{z} $ has simple pole at $0$ and equals $0$ at $\infty $.
\end{example}

\begin{definition}
	We say $R$ is analytic at $\infty $ if $R(\frac{1}{z})$ is analytic at $z=0$.
\end{definition}

\begin{example}
	Polynomial with degree $d$ has a pole of multiplicity $d$ at $\infty $.
\end{example}

\begin{example}
	\[
		\frac{z(z-1)^2(z+3)}{(z+1)^2}
	\]
	has 
	\begin{enumerate}
		\item pole of multiplicity $2$ at $-1$ 
		\item simple roots at $0, -3$ 
		\item double root at $1$ 
		\item double pole at $\infty $
	\end{enumerate}
\end{example}
